\section{Περιγραφή κώδικα ανάλυσης δεδομένων}

Ο κώδικας της ανάλυσης μπορεί έχει χωριστεί σε 4 βασικά κομμάτια,την αίτηση και απόκτηση των δεδομένων από τον server,
την επεξεργασία τους, την άντληση πληροφορίων μέσω οπτικοποίησης και την δημιουργία μηχανικών μοντέλων.


Τα δύο πρώτα βρίσκεται σε ξεχωριστή σελίδα notebook,και για τα πρώτα 2 κόμματια,τα δεδομένα αποθηκεύονται σαν αρχεία σε ειδικό φάκελο.
Το τρίτο και το τέταρτο κομμάτι είναι διαφορετικό πεδία ανάλυσης,αλλά τα έχουμε συγκεντρωμένα σε ένα notebook.
Μπορούν να βρεθούν πολλά διαφορετικά notebook στον αντίστοιχο φάκελο του πηγαίου κώδικα που
σχετίζεται με την αντλήση πληροφορίων μέσω οπτικοποίησης και την δημιουργία μηχανικών μοντέλων και έχουν δοκιμάστει
αρκετές διαφορετικές προσεγγίσεις.
Εμείς θα μιλήσουμε για την τελική προσέγγιση που καταλήξαμε.

\subsection{Απόκτηση δεδομένων από το server}

Στο notebook \texttt{\detokenize{grabTheDataYouWant.ipynb}},χρησιμοποιούμε την βιβλιοθήκη requests για να αντλήσουμε
δεδομένα από τα collection της βάσης δεδομένων,\texttt{userInput} και \texttt{timeSeries}.

Πρώτα παίρνουμε όλα τα δεδομένα του collection \texttt{userInput}
και τα αποθηκεύουμε σε ένα dataframe,με όνομα μεταβλητής \texttt{df\_userInputData}.

Μια σημαντική διευκρίνηση που πρέπει να γίνει είναι ότι το ονόμα αυτό μπορεί να αλλάζει ανά τα διάφορα notebook 
και σε διάφορα σημεία μέσα στο ίδιο notebook,
αλλά πάντα μένει σταθερό ότι οι μεταβλητές που έχουν dataframe και περιέχουν στο όνομα τους \texttt{userInputData},
είναι τα δεδομένα που προέρχονται από το collection \texttt{userInput}.

Συνεχίζουμε ξεδιπλώνοντας σε ξεχωριστές στήλες τα πεδία του \texttt{details} 
και έπειτα επιλέγουμε ποια από τα δεδομένα των χρηστών θα κρατήσουμε με βάση τις τιμές 
που έχει η κεντρική κατάσταση του κάθε πειράματος,του \texttt{experimentState=InsertingSourcePollutant}.
Μαζί με αυτό θα κρατήσουμε την πρώτη και τρίτη κατάσταση του κάθε πειράματος.

Θα κρατήσουμε τα περισσότερα πειράματα που έχουν πραγματοποιηθεί στο δωμάτιο
που ονομάζουμε \textbf{Κρεβατοκάμαρα},αλλά για την μετέπειτα οπτικοποίηση και μηχανική μάθηση,
θα χρησιμοποιήσουμε μόνο τα δεδομένα που προέρχονται από την στήλη του δωματίου με τιμή
\textbf{Κρεβατοκάμαρα id:0 Μ:0.5 Α:1.4 ,id:1 Π:1.7 Δ:0.5  ,id:2 Π:0.4 Α:1.0},που υποδηλώνει τις θέσεις των 
συσκευών στο δωμάτιο,όπως περιγράψαμε στο κεφάλαιο \autoref{sub:Περιγραφή πειραμάτων}.

Από τα επιλεγμένα πειράματα,παίρνουμε την χρονική στιγμή έναρξης του πιο παλιού πειράματος
και την χρονική στιγμή λήξης του πιο καινούργιου πειράματος και θα ζητήσουμε όλα τα δεδομένα των
αισθητήρων που μπορούν να βρεθούν στην βάση δεδομένων μεταξύ αυτών των δύο χρονικών στιγμών στο collection \texttt{timeSeries}.

Αυτή είναι μια ενέργεια που παίρνει πολύ χρόνο και είναι αρκετά σπάταλη,διότι είναι μια χρονική διάρκεια που 
πολύ εύκολα μπορεί να είναι αρκετών εβδομάδων ή μηνών,και θα πετάξουμε την πλειοψηφία αυτών των δεδομένων,
καθώς θα κρατήσουμε μόνο τα δεδομένα που είναι μεταξύ των πειραμάτων που έχουμε επιλέξει.

Αυτό συμβαίνει γιατί χρησιμοποιήσαμε την ίδια αίτηση που κάνουμε και στην οπτικοποίηση της Grafana,
που χρειαζόμαστε όλα τα δεδομένα μεταξύ δύο χρονικών στιγμών.
Αυτό είναι πολύ καλό παράδειγμα απόφασης προγραμματισμού που είναι μια εύκολη απόφαση να πάρεις όταν είσαι στην αρχή,
αλλά πολύ εύκολα μπορεί να ξεφύγει εκτός ελέγχου.
Δεν δημιούργησε μεγάλο πρόβλημα αυτό,καθώς αποθηκεύουμε τα δεδομένα σε αρχείο και δεν ξαναζητάμε δεδομένα από τον server.
Αλλά είναι κάτι κάνει τις αλλαγές στο συγκεκριμένο αρχείο πιο δύσκολες.

Αφού παίρνουμε επιτέλους τα δεδομένα,τα κρατάμε και αυτά σε μορφή dataframe,με το σταθερό όνομα τους 
να είναι το \texttt{timeSeries}.
Διατήρουμε από τα πεδία των δειγματοληψιών των αισθητήρων μόνο την χρονική στιγμή,
το Id του αισθητήρα και το πεδίο σχετικό με την ποιότητα αέρα που έχουμε επιλέξει να χρησιμοποιήσουμε,τα VOC.
Αποθηκεύουμε και το \texttt{timeSeries} και το \texttt{userInputData} σε μορφή json.

Λόγως της δομής των cells των notebooks,προσπαθούμε να μην φτιάξουμε κλάσσεις που να δηλωνόνται στην αρχή και να έχουν 
εκεί όλοι τους την λειτουργικότητα,αλλά να είναι όσο γίνεται πιο αυτόνομα τα κελιά μεταξύ τους και αυτόνομα.
Βεβαίως,Έχουμε φτιάξει σε αρκετές περιπτώσεις συναρτήσησεις,για να μην έχουμε επαναλαμβανώμενο κώδικα,
κυρίως στο κομμάτι της οπτικοποίησης και μηχανικής μάθησης.
Το δύο κύρια αντικείμενο που θα έχουμε καθόλη την διάρκεια της επεξεργασίας και ανάλυσης είναι τα
\texttt{userInputData} και το \texttt{timeSeries}.

\subsection{Επεξεργασία των δεδομένων}

Η επεξεργασία των δεδομένων τελέστηκε κυρίως στο αρχείο \texttt{pre processing data.ipynb},
αλλά κάποια κομμάτια επεξεργασίας έγιναν και στην αρχή του τρίτου notebook,
του \texttt{search data close to sensors.ipynb},αλλά θα ασχοληθούμε με αυτά πιο μετά.
Περιγράφουμε τις βασικές διεργασίες και όχι κάποια αναγκαία φιλτραρίσματα που πρέπει να γίνουν ενδιάμεσα,
όπως είναι να πετάξουμε διπλότυπα (duplicate) δεδομένα ή πειράματα που είναι κενά σε δεδομένα δειγμάτων.


Αφού πάρουμε τα δεδομένα των \texttt{userInputData} και \texttt{timeSeries} από τα json αποθηκευμένα αρχεία 
και τα βάλουμε σε dataframes,συγχονεύουμε τις τρεις καταστάσεις του κάθε πειράματος του \texttt{userInputData}  σε μία,κρατώντας τις χρονικές στιγμές.
Κρατούμε από τα δεδομένα του \texttt{timeSeries} τα δεδομένα που ανοίκουν σε κάποιο πείραμα που έχουμε διατηρήσει,
κρατώντας τα δεδομένα κάθε πειράματος σε ξεχωριστό dataframe για το παρών χρόνο.

Εώς αυτή την στιγμή,ο κάθε αισθητήρας έχει δεδομένα με συχνότητα 3 δευτερόλεπτων ανά δείγμα,
που δεν είναι στις περισσότερες φορές στα ίδια δευτερόλεπτα.
Για να μπορούν να συγκριθούν μεταξύ τους τα δείγματα,επεκτείνουμε ορίζουμε κοινή ελάχιστη και μέγιστη χρονική περίοδο 
για το πείραμα που είναι κοινό για τους τρεις αισθητήρες,ωστέ να έχουν κοινή χρονική διάρκεια.
Έπειτα,για κάθε αισθητήρα,ορίζομυε τα δείγματα τους να αναγράφονται άνα συχνότητα ενός δευτερολέπτου.
Κάνοντας cubic interpolation,μέσω συνάρτησης της pandas, για να συμπληρώσουμε τις ενδιάμεσες στιγμές 
μεταξύ των πραγματικών τιμών.
Τώρα,σε κάθε πείραμα,όλοι οι αισθητήρες έχουν κοινό χρονικό τέλος και αρχή,με όλα τα δεδομένα να είναι ανά δευτερόλεπτο.

Διαμορφώνουμε τα πεδία που εκπροσωπούν τις χρονικές στιγμές,τα \texttt{timestamp},να μην εκφράζουν όλη την χρονική στιγμή,
δηλαδή ώρα και ημερομηνία,αλλά μόνο πόσα δευτερόλεπτα έχουν περάσει από την είσοδο της πηγής.
Με αρνητική τιμή κρατάμε τις τιμές πριν την είσοδο της πηγής στο συγκεκριμένο πείραμα,με θετική τιμή της μετέπειτα τιμές.
Κρατάμε τις αρχικές μορφές των χρονικών διαστημάτων και κάποιες ενδιάμεσες παραμετροποιήσεις,κυρίως στην \texttt{userInputData} ,
αλλά θα χειριζόμαστε τα δεδομένα,και τα \texttt{userInputData} και τα \texttt{timeSeries},μέσω των πεδίων \texttt{seconds}.

Από την \texttt{timeSeries},μετατρέπουμε όποιες αρνητικές τιμές στην στήλη \texttt{VOC} υπάρχουν,
που ξέρουμε ότι η μικρότερη τιμή της είναι το μηδέν,στις μικρότερες θετικές τιμές για το πείραμα στο οποίο βρέθηκαν οι μικρότερες τιμές.

Δημιουργούμε δύο καινούργιες στήλες σχετικές με τα VOC για τα \texttt{timeSeries} κάθε πειράματος.
Η μία στήλη είναι τα VOC με την ελάχιστη τιμή κάθε πειράματος να βρίσκεται στο μηδέν και τις υπόλοιπες να
έχουν μειωθεί κατάλληλα,για να διατήρηθει η αρχική μορφή των VOC.Έχουν ονομαστεί με το ίδιο πεδίο ονόματος
τα \texttt{VOC},με τα αρχικά δεδομένα να αποθηκεύονται στην στήλη με όνομα  \texttt{VOC original}

Η δεύτερη στήλη δημιουργείται εφαρμόζοντας στην καινούργια στήλη \texttt{VOC} rolling windows μεγέθους 30 δειγμάτων
και πέρνοντας για κάθε δευτερόλεπτο την μέση τιμή του αντίστοιχου του rolling windows.
Η στήλη αυτή ονομάζεται \texttt{VOC rolling average} και την μέθοδο της μέσης τιμής των rolling windows την πραγματοποιήσαμε
για να κάνουμε εξομάλυνση (smoothing) των δεδομένων,ωστέ να μειωθεί ο θόρυβος.

Αυτές οι δύο στήλες χρησιμοποιήθηκαν για πολύ καιρό για την ανάλυση των δεδομένων,
Όμως, ύστερα από σύγκριση με την αρχική μορφή των VOC δεδομένων,αποδείχθηκε ότι είχαν χειρότερα αποτελέσματα.
Οπότε θα χρησιμοποιούμε το \texttt{VOC original} από εδώ και πέρα.
Για την εξομάλυνση θα χρησιμοποιήσουμε το Savitzky-Golay φίλτρο,που θα μιλήσουμε αργότερα για αυτό.

Για να περιορίσουμε τον αριθμό δεδομένων που θα λάβουμε υπόψη από κάθε πείραμα,θα βάλουμε ένα κάτω και πάνω όριο
για το χρονικό διάστημα των δεδομένων που θα λάβουμε υπόψη.
Θα κρατήσουμε για κάθε πείραμα,μέχρι και 240 δευτερόλεπτα,δηλαδή 4 λεπτά, πριν από την εισαγωγή της πηγής ρύπου και
μέχρι και 600 δευτερόλεπτα,10 λεπτά,μετά από την εισαγωγή της πηγής ρύπου.

Επιστρέφοντας στα \texttt{userInputData},για κάθε πείραμα θα μετατρέψουμε την θέση πηγής ρύπου 
από την σχετικά απόσταση από τους τοίχους και θα την ορίσουμε σε απόλυτη τιμή \texttt{x axis} και \texttt{y axis}.
Επίσης,ενώνουμε τα δύο αυτά πεδία σε 
Αυτό κάνουμε και για την θέση κάθε συσκευής αισθητήρα,ανάλογη με την θέση αισθητήρων που έχουμε ορίσει στο συγκεκριμένο πείραμα,
ώστε κάθε πείραμα ξεχωριστά να ξέρει ακριβώς που βρίσκονται οι συσκευές αισθητήρα.
Για κάθε πείραμα υπολογίζουμε την ευκλίδεια απόσταση μεταξύ της πηγής ρύπου και των αισθητήρων.

Τα \texttt{timeSeries} των διαφορετικών πειραμάτων τα έχουμε εώς τώρα σε ξεχωριστά dataframes.
Τα ενώνουμε σε ένα μεγάλο dataframe με σκόπο την αποθηκεύση τους σε αρχείο,όπου από εδώ και πέρα θα 
ξεχωρίζουμε μεταξύ τους μέσω της τιμής που θα έχουν στην στήλη \texttt{keys},οπού θα αντιστοιχεί
στο index της γραμμής των λεπτομερειών του πειράματος που εμπεριέχει η \texttt{userInputData}.
Για έξτρα ξεκαθάρηση,το dataframe των \texttt{userInputData},κάθε γραμμή αναφέρεται στις λεπτομέρειες ενός πειράματος
και το dataframe των \texttt{timeSeries} εμπεριέχει τις χρονοσειρές
Τέλος,αποθηκεύουμε τα δύο τελικά dataframe σε καινούργια αρχεία,με μορφή json.

Να συμπληρώσουμε εδώ ότι μετατρέψαμε τις θέσεων των συσκευών των πειραμάτων με 
\texttt{\detokenize{room = Κρεβατοκάμαρα id:1 Π0.6 Α1.2 id:2 Μ0.8 Α1.1}}, που πραγματοποιθήκαν το καλοκαίρι με τους 
παλιούς αισθητήρες BME680,που αντικαταστήσαμε αργότερα.
Οι θέσεις των συσκευών αισθητήρων με id=1 και id=2 των πειραμάτων αυτών,
είναι πολύ κοντά στις συσκευές αισθητήρων των πειραμάτων που συζητάμε, με id=2 και id=0.

Δεν τα χρησιμοποιήσαμε στα μηχανικά μοντέλα εύρεσης πηγής,γιατί η μέθοδος που χρησιμοποιήσαμε 
εύρεσης θέσης της πηγής χρειάζεται τουλάχιστον 3 αισθητήρες,πέρα από ότι ήταν άλλοι αισθητήρες BME680.
Παρόλα αυτά,αυτή η αλλαγή τα κάνει συμβατά για να ενταχθούν σε μελλοντικές μεθόδους,
 με τα πειράματα που έχουμε επιλέξει να εντάξουμε στο τώρα,αν ενσωματόσουμε μεθόδους που να δουλεύουν με λιγότερους
 από τρείς αισθητήρες.


 

Για την επεξεργασία που έγινε στο τελευταίο notebook,το \texttt{search data close to sensors.ipynb},
το κύριο κομμάτι επεξεργασίας είναι να επαναφέρουμε αρκετές από τις στήλες στους τύπους δεδομένων που είχαν πριν την μετατροπή του σε json αρχείο,
όπως και να μεταμορφώσουμε κάποιες στήλες σε μορφές που είναι απαραίτητη για την μετέπειτα ανάπτυξη των μηχανικών μοντέλων.
Ακόμα,στρογκυλοποιούμε τις ευκλίδειες αποστάσεις μεταξύ πηγής και των συσκευών αισθητήρων 

Εδώ είναι που γίνεται εξαμόλυνση του θορύβου των \texttt{VOC original} δεδομένων με το Savitzky-Golay φίλτρο.
Αποτελεί συνάρτηση της βιβλιοθήκης scipy σχετική με τα σήματα (scipy.signal).

Σε αντίθεση με τα rolling windows,που διαμορφώνει κάθε σημείο με βάση των μέσο όρο των γειτονικών δεδομένων,
το Savitzky\-Golay φίλτρο,που η μέθοδος στην scipy ονομάζεται \texttt{savgol\_filter},
κάθε σημείο διαμορφώνεται μετά από την πραγματοποίηση συνέλιξης μεταξύ ενός αριθμού γειτονικών δεδομένων 
και ενός πολυωνύμου βαθμού N.

Επιλέξαμε σαν μέγεθος των γειτονικών δεδομένων,δηλαδή του παραθύρου,να είναι 60 και ο βαθμός πολυωνύμου να είναι το 2.
Έκανε καλύτερη εξαμόλυνση του θορύβου σε σχέση με τα rolling windows,διατηρώντας τα πιο χαρακτηριστικά
τοπικά ελάχιστα και τοπικά μέγιστα των σημάτων του κάθε πειράματος.

Το μόνο κακό αυτήν της μετατροπής δεδομένων,πουδεν αλλάζει την γενική μορφή των δεδομένων,είναι ότι σε σημεία
πολύ ξαφνικών μεγάλων αλλαγών,οι τιμές που βρίσκονται στα τοπικό ελάχιστα κατεβαίνουν σε αρνητικούς αριθμούς,
κάτι που έχουμε καταστήσει ότι κανονικά δεν πρέπει να γίνεται.
Επείδη τα σημεία αυτά είναι μικρά σε εύρος,σπάνια, δεν διαφθήρουν πολύ την κανονική μορφή των δεδομένων,καθώς συμφωνούν με τα υπολοιπα δεδομένα και δεν χρησιμοποιούμε
κάποια μέθοδο ή αλγόριθμο που χρειάζεται καθαρά θετικές τιμές,δεν μας ενοχλεί τόσο.

Ακόμα,υπολογίσαμε για κάθε πείραμα,ποια θα ήταν η ιδεατή στιγμή αυξήσης των \texttt{VOC original} για κάθε αισθητήρα,
με βάση την ευκλίδια απόσταση του αισθητήρα από την πηγή ρύπου.
Αυτό σχετίζεται με μία προσπάθεια προσέγγισης του ιδεατού μοντέλο με το οποίο θεωρήσαμε ότι θα δουλεύψουν
τα δεδομένα και οι τιμές των αισθητήρων.Θα το αναλύσουμε πιο μετά όταν μιλήσουμε για τα αποτελέσματα των μηχανικών μοντέλων.


\subsection{Η άντληση πληροφορίων από την οπτικοποίηση}

Θα κάνουμε κάποια σχόλια πρώτα για το πως αναπτύχθηκε η οπτικοποίηση των δεδομένων και έπειτα
θα κάνουμε σχόλια για αποτελέσματα που φάνηκαν.
Επείδη ο αριθμός των σχεδιαγραμμάτων των δεδομένων είναι αρκετά μεγάλος,δεν θα τις προσθέσουμε σε αυτό το κεφάλαιο.

Η βασική συνάρτηση που βρίσκεται πίσω από την δημιουργία σχεδιαγραμμάτων των 
\texttt{timeSeries} του κάθε πειράματος,
είναι η \texttt{printDataBasedOnDate}.

Σε αντίθεση με ότι υπονοεί το όνομα της συνάρτησης για σχεδιασμό μέσω ημερομηνίας,
Η συνάρτηση αυτή σχεδιάζει τα \texttt{timeSeries} δεδομένα των πειράματων,ομαδοποιημένα με βάση
τις διακριτές τιμές οποιαδήποτε στήλη της \texttt{userInputData} επιθυμούμε,δίνοντας το αντίστοιχο όνομα της στήλης.

Όταν μιλάμε ομαδοποίηση με διακριτές τιμές,εννοούμε ότι βάζουμε σε μία κοινή ομάδα όσα 
πειράματα έχουν την ίδια τιμή για το πεδίο,για παράδειγμα όλα τα πειράματα που γίνουν στις
14 Οκτώμβρη,και το κάνουμε αυτό για όλες τις τιμές της στήλης που μας δώθηκε που είναι ξεχωριστές.
Χρησιμοποιούμε τον όρο διακριτές και όχι ξεχωριστές,επειδή δεν είναι ότι έχουν κάποια ξεχωριστή 
ιδιότητα,αλλά ότι είναι διαφορετικές μεταξύ τους.

Αν τις δώσουμε πολλαπλά ονόματα στηλών της \texttt{userInputData},
θα ομαδοποιηθούν με σειρά αντίθετης της σειράς που τα εισάγαμε.
Δηλαδή,θα ξεκινήσουμε την επανάληψη ομαδοποιήσεων αλλάζοντας τις διακριτές τιμές τις τελευταίας στήλης,
δηλαδή της στήλης που δωθήκε πιο δεξιό στην λίστα ονομάτων που δώσαμε στις παραμέτρους της συνάρτησης,
κρατόντας τις άλλες στήλες σταθερές και βλέποντας ποια πειράματα υπάρχουν για αυτές τις διακριτές τιμές
των στηλών της \texttt{userInputData}
Αφού εξαντληθούν όλες οι διακριτές τιμές της τελευραίας στήλης,
θα αλλάξουμε την διακριτή τιμή της προτελευταίας στήλης,το όνομα που είναι δεύτερο πιο δεξία, και θα ξαναεπαναλάβουμε τις διακρητές τιμές της τελευταίας στήλης,
και πάει λέγοντας μέχρι να τελειώσουν οι διακριτές τιμές της προτελευταίας στήλης,οπού μετά πάμε στην αμέσως επόμενη διακριτή τιμή.
Όλο αυτό συνεχίζει μέχρι να δούμε όλες τις δυνατές ομαδοποιήσεις των στηλών που μας δώθηκαν.

Ο λόγος που εστιάζουμε τόσο στην ομαδοποίηση με διακριτές τιμές είναι διότι αποτελεί ένα πολύ
δυνατό εργαλείο διαχήρισης των dataframes της pandas, 
την μέθοδο \texttt{groupby()},το οποίο χρησιμοποιούμε πολύ από εδώ και πέρα ,
για να διαχειριστούμε  κυρίως τα \texttt{timeSeries},αλλά και τα \texttt{userInputData}.

Κρατάμε σταθερά τα χαρακτηριστικά χρώματα που είχαμε και στην grafana,δηλαδή μπλε για το αισθητήρα
με id=0,πράσινο για id=1 και κίτρινο για id=2